\newpage
\chapter*{\centering{\MakeUppercase{Lời nói đầu}}}
\addcontentsline{toc}{chapter}{\textcolor{fitblue}{\MakeUppercase{Lời nói đầu}}}

\noindent \indent Đồ án Caro được xây dựng trong khuôn khổ môn học Cơ sở lập trình tại Trường Đại học Khoa học Tự nhiên TP.HCM. Thay vì chỉ dừng lại ở các giao diện console cơ bản, dự án của nhóm đã tập trung ứng dụng ngôn ngữ C++17 kết hợp cùng Win32 API và thư viện đồ họa GDI+ để tạo ra một thế giới bóng đá Pixel Art sinh động.

Đồ án không chỉ là một trò chơi trí tuệ, mà còn là sự kết hợp độc đáo giữa logic cờ caro và không khí cuồng nhiệt của các sân vận động hàng đầu thế giới. Dự án tích hợp nhiều tính năng hiện đại như:
\begin{itemize}
    \item Hệ thống AI thông minh: Sử dụng thuật toán Alpha-Beta Pruning và Zobrist Hashing để tạo ra các cấp độ thử thách từ cơ bản đến chuyên nghiệp.

    \item Giao diện \& Âm thanh sống động: Hiệu ứng Glassmorphism hiện đại, hệ thống âm thanh đa luồng không gây gián đoạn và các nhân vật biểu tượng như Ronaldo, Messi, Neymar dưới dạng Pixel Art.

    \item Cơ chế chuyên sâu: Thống kê tỷ lệ kiểm soát bóng, số nước đi, hệ thống lưu/tải trò chơi nhị phân đảm bảo tính ổn định và bảo mật dữ liệu.
\end{itemize}

Thông qua đồ án này, nhóm chúng em không chỉ củng cố kỹ thuật xử lý mảng, tập tin và con trỏ mà còn học cách tổ chức mã nguồn theo mô-đun chuyên nghiệp và tối ưu hóa hiệu suất ứng dụng.

Chúng em xin gửi lời cảm ơn chân thành đến thầy Trương Toàn Thịnh đã tận tình hướng dẫn và tạo điều kiện để nhóm có thể hoàn thiện đồ án này một cách tốt nhất.
